\chapter{Acknowledgements}

% Acknowledgments\ldots.

This book owes much to two German scholars who are sadly not with us any more. I am sorry for not being able to present this book in front of them. They are the late Professor Dr.\ Wilhelm Halbfass, and the late Professor Dr.\ Tillman Vetter. 

Professor Halbfass was the patient supervisor of my doctoral thesis in my five and a half years at Penn in Philadelphia. Knowing my interest in the \PYSV, he almost matter-of-factly suggested that I prepare a critical edition of the text as the topic for my doctoral dissertation. I did not have to think much before saying yes. He also made it possible that I had access to manuscripts. The thesis became the basis of this book. I am not happy that I became one of his very last students. 

Professor Vetter practically determined the content of this book. Even before I met him, a brief remark in one of his writings directed me to look at the Īśvara section of the Vivaraṇa. I was fortunate to meet him years later, as a Gonda Fellow in Leiden in 2005 to 2006. He agreed to read my edition with me. In view of the limited time period, I asked him to read the Īśvara section in which I thought he might be interested. In the end, it is not a coincidence that the content of this book corresponds to the part we read together, including the commentary on sūtra 1.1. It was also he who insisted that I have a translation to read the edition with him. 

I also owe much to Professor Albrecht Wezler. He agreed that I work on a critical edition of the \PYSV\ and let me use his copies of the Lahore manuscript which has always been difficult to gain access to. After I came to Hamburg in 2000, he arranged weekly reading sessions for the edition I prepared for my thesis. Much of the improvements in the text come from those reading sessions between 2000 and 2002.

I am thankful to all my teachers. Professor Osamu Hayashima first taught me Sanskrit. It might be that I am the only person who still reads Sanskrit among his students. The way he taught me Sanskrit was very efficient; he told me to learn it by myself. Professor Hiromasa Tosaki imparted to us what it is to be a scholar. Professor Akihiko Akamatsu introduced me to the concept of philology. It took some time for me to digest the concept but I hope now I get it. I have learned many many things from Dr.\ Futoshi Omae when he was an assistant at the Kyushu University. It might have been during the daily chatting over many things with him that I learned most about our field of studies.

When I was at Penn, it was a golden age for classical indological studies there. We had Professors Halbfass, George Cardona, Ludo and Rosane Rocher and many inspiring fellow students. I learned various things from all of them. I have so many fond memories of my time in Philadelphia. One of them is the memory of reading the Vākyapadīya with Prof.\ Cardona for five years.

Working for the Skandapurāṇa project, led by Professor Hans Bakker, in Groningen, the Netherlands, was a new experience for me. Exposure to the Dutch Indology had an effect on me. I had to reshape my self-perception as an indologist. I also thank Prof.\ Bakker's support/insistence to finish this book as soon as possible. %At the same time, I feel apologetic that it took so long.

The same thanks go to Arlo Griffiths whom I first met in 1997 or so at an AOS meeting before seeing each other again in Groningen in 2004. He, too, has been supportive and persuasive. He arranged that I could read my text with Professor Vetter. He also looked at my English many times. I was also fortunate to have had an opportunity to work with Professor Yuko Yokochi in Groningen. In 2012, she arranged an intensive course to read the Vivaraṇa at the Kyoto University. Inputs from scholars and students there were very valuable. The people who were present there include Diwakar Acharya, Somdev Vasudeva, Akihiro Kanabishi, Koreto Ikehata. 

Special mention has to be made to Harunaga Isaacson with whom I seem to share some fate in that we have been moving symmetrically around Hamburg in the past 14 years. Not only has he been a tremendous inspiration, without him, this book would have been impossible. If there is any quality in it, it is largely because of him. I sincerely thank him for his support, not only for this book but also for everything. 

There are many other people to whom thanks are due. Here are some of the names: 
% Without any of the people mentioned below, this work would have been impossible. It took too long for any standard to get this out. I try here to mention those whom this book owes in generally a chronological order.
% Osamu Hayashima
% Hiromasa Tosaki
% Akihiko Akamatsu
% Wilhelm Halbfass
% George Cardona, 
% Ludo Rocher 
% Rosane Rocher
% Hans Bakker
% Albrecht Wezler
% Harunaga Isaacson
% Arlo Griffiths
% Diwakar Acharya
% Friends at PENN
% Friends at RuG
% Friends in Fukuoka
% Friends in Hamburg
% Korean scholars
% Tillman Vetter
% Yuko Yokochi
% Diwakar Acharya,
Ashok Aklujkar,
Orna Almogi,
Loriliai Biernacki,
Piotr Balcerowicz,
Benita von Behr,
Bidur Bhattarai,
Peter Bisschop,
Tim Cahill,
Britton Chance\dag{},
Danielle Cuneo,
Jonathan Duquette,
Christina Edingloh,
John Freeman,
Camillo A. Formigatti,
Angelika Frankowski,
Elisa Freschi,
Takamichi Fujii,
Dominic Goodall,
Kunihiro Harikai,
Tomoko Himeno,
%Koreto Ikehata,
%Akihiro Kanabishi,
Kazuo Kano,
Kei Kataoka,
Kenichi Kuranishi,
Andrey Klebanov,
Masato Kobayashi,
Stasinos Konstantopoulos,
Sonja Lakner,
Philipp Maas,
Kristin McGee,
David Mellins,
Mai Moriguchi,
Shin'ya Moriyama,
Yasutaka Muroya,
Hodo Nakamura,
Masanori Nakagawa,
David Nelson, 
Heidi Neubert,
Sachio Nioka,
Shoko Nioka,
Luther Obrock,
Kiyokazu Okita,
Dimitri Pauls,
Kim Plofker,
Akane Saito,
Florinda de Simini,
Barbara Schuler,
Greg Seton,
Francesco Sferra,
Taisei Shida,
Kiyokuni Shiga,
Jae\-kwan Shim,
Iain Sinclair,
Michael Slouber,
Hans-Jörg Stein,
Vidyasankar Sundaresan,
Peter Szanto,
Masaya Takahashi,
Ryugen Tanemura,
Somdev Vasudev,
Dorji Wangchuk,
Eiichi Yamaguchi%
.


Since it took so long to complete this book, practically anyone I knew has been in one way or another influential. I would like to thank all the friends I have ever had. And at last but not least, the biggest thanks are due to my parents for their unwavering support.
% Koh Endo
% Professor Osamu Hayashima first taught me Sanskrit. For a person like myself, his method of teaching Sanskrit turned out to be best; he was very strict. Professor Hiromasa Tosaki taught me what it means to be a researcher. Without Dr.\ Futoshi Omae, I would not have worked on this intriguing text, the Pātañjalayogaśāstravivaraṇa. I still remember the moment he, then the assistant of the department and history of Indian philosophy and buddhism, Kyushu University, asked me what I wanted to work on for my master's thesis. When I mentioned the Yoga school, he entered the library and came back with the edition of ``the Vivaraṇa.'' He said that I could work on it for ten years; it has been more than twenty years and I am not done yet. I owe much to Professor Akihiko Akamatsu for being patient and not giving up on me. I cannot help feel a sense of irony because the topic of this book is essentially the same as the master's thesis I wrote with him. After more than twenty years, I would hope that I learned something more and this book offers something new.

% To this day, I regret the loss of my dissertation advisor, Professor Wilhelm Halbfass. The inspirations he gave me were enormous; his knowledge and intellect were aspiring; and he was kind. He helped me through obtaining the degree at the University of Pennsylvania. This work, nonetheless, is based on my doctoral dissertation. It was one of my happiest moments when he said, ``Now you have convinced me,'' with regard to the authorship when I showed the final draft of my dissertation. It is a pity that he passed away only five months after I had my defense. He is not there any more to testify that he said that.

% I also owe much to Professors George Cardona, Ludo Rocher and Rosane Rocher of the University of Pennsylvania. Having Prof.\ Cardona as a teacher is a great advantage. He makes it clear that he never openly attacks his student (however he may be wrong). It is a great honor that he considers myself as one of his students. 